\documentclass[11pt]{article}
\usepackage{amsmath,amssymb,amsthm}
\usepackage{geometry}
\geometry{margin=1in}

\title{Metoda liczenia liczb Ramsey na przykładzie \(R(5,5)=43\): Shadow Analysis of Critical Cores}
\author{Michał Machura}
\date{}

\begin{document}
\maketitle


% BEGIN_MACHURA_SHADOW_OBSERVATIONS
\section{Observations from shadow-threshold plots}

The starting point of the method was not a direct search for a larger Ramsey graph, but the repeated observation of stable boundary obstructions in computational experiments. In the case of \(R(5,5)\), computations at the boundary layer \(42 \to 43\) repeatedly produced a minimal obstruction value equal to \(2\). This behavior was compared with smaller Ramsey families, where the boundary obstruction takes different values.

The key observation is that the boundary shadow value is not a universal constant. It depends on the Ramsey family \(R(a,b)\), the critical-core layer, and the corresponding one-vertex-extension obstruction.

\paragraph{Remark.}
The quantity
\[
\Delta_{a,b}(k \to k+1)
\]
is specific to each Ramsey number \(R(a,b)\). There is no universal constant. This explains why a direct method based only on searching for graphs or minimizing an unrestricted score is insufficient. The correct object is the shadow-boundary value associated with clean critical cores.

For the cases already recorded in the present calculation, the observed shadow-boundary values include:
\[
\Delta_{3,3}(5 \to 6)=2,
\]
\[
\Delta_{3,4}(8 \to 9)=2,
\]
\[
\Delta_{3,5}(13 \to 14)=4,
\]
\[
\Delta_{3,6}(17 \to 18)=4,
\]
\[
\Delta_{3,7}(22 \to 23)=4,
\]
\[
\Delta_{4,4}(17 \to 18)=12,
\]
and for the main case:
\[
\Delta_{5,5}(42 \to 43)=2.
\]

Thus, the value \(2\) in the \(R(5,5)\) calculation is not an arbitrary search artifact. It is the shadow-boundary obstruction value for the \(42 \to 43\) layer.

The plots used during the discovery phase showed three important features:

\begin{enumerate}
    \item For \(R(5,5)\), the boundary layer at \(n=43\) stabilizes at obstruction value \(2\).
    \item Other Ramsey families exhibit different boundary values, for example \(4\) or \(12\).
    \item The edge-complement symmetry and the red/blue conflict split support the interpretation of a boundary layer rather than a random search failure.
\end{enumerate}

Consequently, the method records
\[
shadow\_boundary\_delta(R(a,b))
\]
as the main table value, while unrestricted low-score graph searches are treated only as diagnostic information.
% END_MACHURA_SHADOW_OBSERVATIONS

\end{document}

